\section{Study texts and liturgical locators}
The biblical English in this section is the public-domain Douay--Rheims/Challoner translation. It is \textbf{not the approved English proclaimed at Mass}; the U.S. Lectionary controls proclamation. Full verse envelopes used for context are labelled below. Missal adaptations, excerpts and prayers are identified by their own locators rather than reconstructed from a Bible translation.

\subsection{Entrance}
\label{proper-entrance}
\textit{Ego clamavi}. Roman Missal, Week XXIX, Entrance; cf. Psalm 16:6,8 (Vulgate), 17:6,8 (Hebrew). The Missal's adapted antiphon text is not reproduced. Its two scriptural images are the cry for hearing and divine shelter; the full Psalm supplies the context discussed above.

\subsection{Collect}
\label{proper-collect}
\textit{Omnipotens sempiterne Deus, fac nos}. Roman Missal, Week XXIX, Collect. Approved English is not reproduced; the commentary explains the will's disposition and sincere service without supplying a substitute prayer.

\subsection{First reading}
\label{proper-first-reading}
Isaiah 45:1,4--6; U.S. Lectionary no.145. Verses 2--3 are not in this appointment.
\begin{englishwitness}{Douay--Rheims/Challoner study English}
\textsuperscript{1} Thus saith the Lord to my anointed Cyrus, whose right hand I have taken hold of, to subdue nations before his face, and to turn the backs of kings, and to open the doors before him, and the gates shall not be shut.

\textsuperscript{4} For the sake of my servant Jacob, and Israel my elect, I have even called thee by thy name: I have made a likeness of thee, and thou hast not known me.

\textsuperscript{5} I am the Lord, and there is none else: there is no God besides me: I girded thee, and thou hast not known me:

\textsuperscript{6} That they may know who are from the rising of the sun, and they who are from the west, that there is none besides me. I am the Lord, and there is none else:
\end{englishwitness}

\subsection{Responsorial Psalm}
\label{proper-psalm}
Psalm 96:1,3,4--5,7--8,9--10a,c (Hebrew; Vulgate/Douay 95), response 7b; Lectionary no.145. The response uses part of verse 7, not the whole verse. The following study text gives complete verses 7 and 10: verse 10's middle clause about the world's stability is \textbf{context only and is not appointed}. It is retained here visibly as part of a complete Bible verse, not restored to the sung strophe.
\begin{englishwitness}{Douay--Rheims/Challoner whole-verse study envelope}
\textsuperscript{1} A canticle for David himself, when the house was built after the captivity. Sing ye to the Lord a new canticle: sing to the Lord, all the earth.

\textsuperscript{3} Declare his glory among the Gentiles: his wonders among all people.

\textsuperscript{4} For the Lord is great, and exceedingly to be praised: he is to be feared above all gods.

\textsuperscript{5} For all the gods of the Gentiles are devils: but the Lord made the heavens.

\textsuperscript{7} Bring ye to the Lord, O ye kindreds of the Gentiles, bring ye to the Lord glory and honour:

\textsuperscript{8} Bring to the Lord glory unto his name. Bring up sacrifices, and come into his courts:

\textsuperscript{9} Adore ye the Lord in his holy court. Let all the earth be moved at his presence.

\textsuperscript{10} Say ye among the Gentiles, the Lord hath reigned. For he hath corrected the world, which shall not be moved: he will judge the people with justice.
\end{englishwitness}

\subsection{Second reading}
\label{proper-second-reading}
1 Thessalonians 1:1--5b; Lectionary no.145. The reading ends before verse 5's final clause about the missionaries' conduct. The complete verse 5 below includes that final clause as \textbf{unappointed context}.
\begin{englishwitness}{Douay--Rheims/Challoner whole-verse study envelope}
\textsuperscript{1} Paul and Sylvanus and Timothy to the church of the Thessalonians: in God the Father and in the Lord Jesus Christ.

\textsuperscript{2} Grace be to you and peace. We give thanks to God always for you all: making a remembrance of you in our prayers without ceasing,

\textsuperscript{3} Being mindful of the work of your faith and labour and charity: and of the enduring of the hope of our Lord Jesus Christ before God and our Father.

\textsuperscript{4} Knowing, brethren, beloved of God, your election:

\textsuperscript{5} For our gospel hath not been unto you in word only, but in power also: and in the Holy Ghost and in much fulness, as you know what manner of men we have been among you for your sakes.
\end{englishwitness}

\subsection{Alleluia verse}
\label{proper-acclamation}
Philippians 2:15d,16a; Lectionary no.145. Only the short portions concerning shining in the world and holding the word of life are appointed. The rest of these complete verses is context, not additional acclamation text.
\begin{englishwitness}{Douay--Rheims/Challoner whole-verse study envelope}
\textsuperscript{15} That you may be blameless and sincere children of God, without reproof, in the midst of a crooked and perverse generation: among whom you shine as lights in the world.

\textsuperscript{16} Holding forth the word of life to my glory in the day of Christ: because I have not run in vain, nor laboured in vain.
\end{englishwitness}

\subsection{Gospel}
\label{proper-gospel}
Matthew 22:15--21; Lectionary no.145. No shorter form is appointed for this Year A Gospel.
\begin{englishwitness}{Douay--Rheims/Challoner study English}
\textsuperscript{15} Then the Pharisees going, consulted among themselves how to insnare him in his speech.

\textsuperscript{16} And they sent to him their disciples with the Herodians, saying: Master, we know that thou art a true speaker and teachest the way of God in truth. Neither carest thou for any man: for thou dost not regard the person of men.

\textsuperscript{17} Tell us therefore what dost thou think? Is it lawful to give tribute to Caesar, or not?

\textsuperscript{18} But Jesus knowing their wickedness, said: Why do you tempt me, ye hypocrites?

\textsuperscript{19} Shew me the coin of the tribute. And they offered him a penny.

\textsuperscript{20} And Jesus saith to them: Whose image and inscription is this?

\textsuperscript{21} They say to him: Caesar’s. Then he saith to them: Render therefore to Caesar the things that are Caesar’s; and to God, the things that are God’s.
\end{englishwitness}

\subsection{Prayer over the Offerings}
\label{proper-offerings}
\textit{Tribue nos, Domine, quaesumus}. Roman Missal, Week XXIX, Prayer over the Offerings. Approved English is not reproduced. The discussion concerns gifts, service and purification, not a replacement translation.

\subsection{Communion: Psalm alternative}
\label{proper-communion-psalm}
\textit{Ecce oculi Domini}. Roman Missal, Week XXIX, first Communion alternative; cf. Psalm 32:18--19 (Vulgate), 33:18--19 (Hebrew). This is an adaptation. The Missal antiphon is not reproduced or re-created by joining English Bible phrases.

\subsection{Communion: Mark alternative}
\label{proper-communion-mark}
\textit{Filius hominis venit}. Roman Missal, Week XXIX, second Communion alternative; Mark 10:45. The appointed excerpt concerns the Son of Man's gift of his life as ransom. The verse's preceding service clause belongs to the Gospel context and is not additional antiphon text. The antiphon is not reproduced. This alternative and the Psalm alternative are not successive obligatory chants.

\subsection{Prayer after Communion}
\label{proper-after-communion}
\textit{Fac nos, quaesumus, Domine}. Roman Missal, Week XXIX, Prayer after Communion. Approved English is not reproduced. The analysis concerns present sacramental benefit and preparation for eternal gifts.
